% Results section — PRD two-column
% REVISED: production run values pending; thesis baselines used as placeholders

Having established the pipeline in Sec.~\ref{sec:method}, we now present
the results of the EMRI dark siren analysis. We begin with the properties
of the simulated detection catalog (Sec.~\ref{sec:detections}), validate
the parameter estimation against known benchmarks
(Sec.~\ref{sec:pe_validation}), and then present the central result:
the $\HH$ posterior with and without the MBH mass channel
(Sec.~\ref{sec:h0_posterior}). We conclude with the impact of
detection probability and catalog completeness corrections
(Sec.~\ref{sec:pdet_completeness}) and a discussion of systematic
effects (Sec.~\ref{sec:systematics}).

%% ============================================================
\subsection{Simulated EMRI detections}
\label{sec:detections}
%% ============================================================

We simulate EMRI events across a grid of seven trial Hubble constant
values $h \in \{0.60, 0.65, 0.70, 0.73, 0.75, 0.80, 0.90\}$, drawing
source parameters from the astrophysical distributions described in
Sec.~\ref{sec:method}. A total of $N_\mathrm{inj} = 165{,}000$ events
are injected, and the SNR is computed for each using the full LISA TDI
response~\cite{Babak:2023lro}. An event is considered detected when its
network $\SNR$ exceeds the threshold $\SNR_\mathrm{thr} = 15$.

Of the $165{,}000$ injected events, $N_\mathrm{det} = 663$ pass the
detection threshold across all $h$-values, corresponding to an overall
detection fraction $f_\mathrm{det} = 4.0 \times 10^{-3}$. The per-$h$
detection yield increases monotonically from
$f_\mathrm{det}(h=0.60) = 2.2 \times 10^{-3}$ to
$f_\mathrm{det}(h=0.90) = 8.1 \times 10^{-3}$, a factor of $3.6$
increase. This trend is expected: larger $\HH$ implies smaller luminosity
distances at fixed redshift, hence higher SNR. Minor non-monotonicity at
$h = 0.70$ is consistent with Poisson noise at the $1\sigma$ level.

All detected events lie at redshift $z < 0.21$, well below the injection
cutoff at $z_\mathrm{cut} = 0.5$, which provides a safety margin of at
least $2.5\times$. The Farr criterion~\cite{Farr:2019rap}
$N_\mathrm{inj}/N_\mathrm{det} \geq 4$ is satisfied for all $h$-values,
with the minimum ratio of $124$ occurring at $h = 0.90$.

For the posterior evaluation, we further require that each detection has a
relative luminosity distance uncertainty
$\sigma_{\dL}/\dL < 0.10$, as estimated from the Cram\'{e}r--Rao
bound. This cut reduces the sample to
\pending{N\_det posterior after $\sigma_{d_L}/d_L < 0.10$ cut from
production run} detections.

\begin{figure}[t]
  \centering
  % \includegraphics[width=\columnwidth]{figures/snr_distribution.pdf}
  \todo{Figure: Left panel --- histogram of $\SNR$ for detected events
  showing the peak near threshold and power-law tail.
  Right panel --- scatter of $\SNR$ vs $\dL$ colored by redshift,
  showing the detection horizon.}
  \caption{Distribution of detected EMRI events. \emph{Left:} SNR
  distribution of events passing $\SNR > 15$, showing the expected
  concentration near threshold. \emph{Right:} SNR versus luminosity
  distance, colored by source redshift. The detection horizon at
  $\dL \approx 1\,\mathrm{Gpc}$ is clearly visible.}
  \label{fig:snr_distribution}
\end{figure}

%% ============================================================
\subsection{Parameter estimation validation}
\label{sec:pe_validation}
%% ============================================================

The Fisher information matrix for each detected EMRI is computed using
forward-difference derivatives of the full LISA TDI waveform
\cite{Vallisneri:2007ev}, yielding Cram\'{e}r--Rao bounds on the
14-dimensional EMRI parameter space. We validate the Fisher matrix
computation against Table~III of Babak
\emph{et~al.}~\cite{Babak:2017tow}, which provides benchmark
Cram\'{e}r--Rao bounds for a fiducial EMRI source. Our results agree with
the benchmark values to within the expected numerical precision of the
finite-difference stencil.

For the dark siren inference, three parameter estimation outputs are
critical: the luminosity distance uncertainty $\sigma_{\dL}$, the sky
localization area $\Delta\Omega$, and the redshifted MBH mass uncertainty
$\sigma_{\Mz}$. The sky localization is computed as
\begin{equation}
  \Delta\Omega = 2\pi |\sin\theta_S| \,
  \sqrt{\sigma_{\theta_S}^2 \, \sigma_{\phi_S}^2
    - \mathrm{cov}(\theta_S, \phi_S)^2} \,,
  \label{eq:sky_localization}
\end{equation}
where $\sigma_{\theta_S}$ and $\sigma_{\phi_S}$ are the
Cram\'{e}r--Rao bounds on the ecliptic co-latitude and longitude,
respectively.

The redshifted MBH mass $\Mz$ is measured with extraordinary precision
by the gravitational waveform: the fractional uncertainty is
$\sigma_{\Mz}/\Mz \lesssim 10^{-6}$ for sources with
$\SNR \geq 15$~\cite{Babak:2017tow}. In contrast, the BH mass of
candidate host galaxies, inferred from the stellar mass via the
$\Mbh$--$\Mstellar$ scaling relation~\cite{Reines:2015moa}, carries a
fractional uncertainty of $\sigma_{\Mbh}/\Mbh \sim 0.1$. This
difference of roughly five orders of magnitude justifies treating
$\Mz$ as exactly known in the likelihood evaluation: the GW mass
measurement acts as an effective delta function when compared against the
broad galaxy mass distribution, i.e.,
\begin{equation}
  p(\Mz^\mathrm{obs} \mid \Mz^\mathrm{true})
  \approx \delta(\Mz^\mathrm{obs} - \Mz^\mathrm{true}) \,.
  \label{eq:delta_mass}
\end{equation}
This approximation eliminates $\Mz$ as a free parameter in the GW
likelihood, so that the mass information enters solely through the
requirement that the candidate host galaxy's BH mass satisfies
$M_\mathrm{bh} = \Mz / (1+z)$.

%% ============================================================
\subsection{$\HH$ posterior: with vs.\ without MBH mass}
\label{sec:h0_posterior}
%% ============================================================

We evaluate the $\HH$ posterior on a grid of $h$-values spanning
$[0.60, 0.90]$ using the dark siren likelihood described in
Sec.~\ref{sec:method}. Figure~\ref{fig:posterior_comparison} shows the
central result of this work: the comparison of the $\HH$ posterior
obtained \emph{without} the MBH mass channel (using $\dL$, $\phi_S$,
$\theta_S$ only) and \emph{with} the MBH mass channel (additionally
using $\Mz$).

\paragraph{Without MBH mass.}
\pending{Rank-1 / production value for $h^\mathrm{(no\,mass)}$}. The
thesis-era baseline ($\pdet = 1$) is withdrawn: it predates the current
pipeline and no surviving result file reproduces it.

\paragraph{With MBH mass.}
\pending{Rank-1 / production value for $h^\mathrm{(mass)}$}. The
thesis-era baseline ($\pdet = 1$) is withdrawn: it predates the current
pipeline and no surviving result file reproduces it.

\paragraph{Mass-channel ablation.}
\pending{Rank-1: $\sigma_h$ with and without the mass term, conditional
and marginalised, per scenario.}

\begin{figure}[t]
  \centering
  % \includegraphics[width=\columnwidth]{figures/h0_posterior_comparison.pdf}
  \todo{Figure: $\HH$ posterior comparison. Blue solid --- without MBH
  mass. Red dashed --- with MBH mass. Vertical gray line at $h = 0.73$
  (injected value). Shaded $1\sigma$ bands.}
  \caption{Posterior on the Hubble constant $\HH$ from
  \pending{N\_det} EMRI dark siren events. The blue curve uses only
  the luminosity distance and sky position; the red curve additionally
  incorporates the redshifted MBH mass. The vertical
  dashed line marks the injected value $h = 0.73$.}
  \label{fig:posterior_comparison}
\end{figure}

\pending{Rank-1: single-event likelihoods $p(d_\mathrm{GW}^i \mid h)$
with and without the mass term (Fig.~\ref{fig:single_event}).}

\begin{figure}[t]
  \centering
  % \includegraphics[width=\columnwidth]{figures/single_event_likelihoods.pdf}
  \todo{Figure: Single-event likelihoods for 4 representative
  detections. Left column: without MBH mass (broad/flat). Right column:
  with MBH mass.}
  \caption{Single-event likelihoods $p(d_\mathrm{GW}^i \mid h)$ for
  four representative EMRI detections at different redshifts.
  \emph{Left column:} using only $\dL$ and sky position --- the
  likelihoods are broad due to the large number of candidate hosts.
  \emph{Right column:} additionally using $\Mz$.}
  \label{fig:single_event}
\end{figure}

\paragraph{Convergence with event number.}
To assess the scaling of the $\HH$ constraint with the number of
detections, we evaluate the posterior on random subsets of
\pending{subset size} events drawn from the full catalog. The
posterior width narrows approximately as $N_\mathrm{det}^{-1/2}$
for both channels, as expected for independent measurements.
\pending{Rank-1: posterior width against the number of events, per
likelihood channel.}

\begin{figure}[t]
  \centering
  % \includegraphics[width=\columnwidth]{figures/posterior_convergence.pdf}
  \todo{Figure: $1\sigma$ width of $\HH$ posterior vs number of events
  for both channels. Dashed line: $N^{-1/2}$ scaling.}
  \caption{Convergence of the $\HH$ posterior width with the number of
  EMRI detections. Points show the median $1\sigma$ width from
  \pending{N\_subsets} random subsets; error bars indicate the
  16th--84th percentile range. The dashed line shows the expected
  $N^{-1/2}$ scaling.}
  \label{fig:posterior_convergence}
\end{figure}

%% ============================================================
\subsection{Detection probability and completeness effects}
\label{sec:pdet_completeness}
%% ============================================================

The selection-corrected results with a realistic $\pdet$ are
\pending{Rank-1 budget}. We now describe the treatment of $\pdet$ and of
the catalog completeness correction.

\paragraph{Detection probability.}
The detection probability $\pdet(\dL, \Mbh, h)$ is estimated via
importance-sampling--weighted kernel density estimation on the injection
campaign (Sec.~\ref{sec:method}), validated with Wilson confidence
interval overlap across $916$ grid bins with zero
Benjamini--Hochberg--adjusted false
discoveries~\cite{Farr:2019rap,Tiwari:2017ndi}.
\pending{Present results with realistic $\pdet$ from production run.}

\paragraph{Catalog completeness.}
The GLADE+ galaxy catalog~\cite{Dalya:2021ewn} is incomplete at redshifts
$z \gtrsim 0.03$. To avoid the systematic bias toward low $\HH$ that
arises when the catalog is denser at lower redshifts (where smaller $h$
maps the observed $\dL$ to a nearer, better-cataloged volume), we
implement the completeness correction of Gray
\emph{et~al.}~\cite{Gray:2019ksv}. The corrected per-event likelihood is
\begin{equation}
  p_i(x_\mathrm{GW} \mid h) =
    f_i \cdot \Lcat^{(i)}
    + (1 - f_i) \cdot \Lcomp^{(i)} \,,
  \label{eq:completeness_correction}
\end{equation}
where $f_i = f(z_i, h) \in [0, 1]$ is the completeness fraction of the
GLADE+ catalog at the inferred redshift $z_i = z(\dL^{(i)}, h)$, and
$\Lcat^{(i)}$ and $\Lcomp^{(i)}$ are the catalog and completion
likelihood terms, respectively. The completion term accounts for the
probability of a host galaxy not present in the catalog:
\begin{equation}
  \Lcomp^{(i)} =
    \frac{\displaystyle\int dz \; \pGW(x \mid z, \Omega_\mathrm{det}, h)
      \; \pdet(\dL(z,h)) \; \frac{dV_c}{dz}}
    {\displaystyle\int dz \; \pdet(\dL(z,h))
      \; \frac{dV_c}{dz}} \,.
  \label{eq:completion_term}
\end{equation}

The completeness fraction $f(z, h)$ is interpolated from the GLADE+
completeness data. At $h = 0.73$, the catalog is approximately $50\%$
complete at $z \approx 0.03$, dropping to $\sim 35\%$ at
$z \approx 0.11$ and $\sim 21\%$ at $z \gtrsim 0.18$, where it
plateaus due to extrapolation beyond the available data. The $h$-dependence
arises because larger $\HH$ maps a given $z$ to a smaller $\dL$, at which
the catalog is more complete: $f(z = 0.05,\, h = 0.86) > f(z = 0.05,\, h = 0.60)$.

The combination formula~\eqref{eq:completeness_correction} satisfies two
important limiting cases: $f_i = 1$ exactly recovers the catalog-only
likelihood, and $f_i = 0$ gives the pure completion term, which uses only
the comoving volume prior and the GW likelihood. These limits, together
with the requirement that the combined likelihood interpolates
monotonically between the two terms, have been verified analytically and
numerically.

\pending{Show the effect of the completeness correction on the $\HH$
posterior bias (production run, $h_\mathrm{true} = 0.73$).}

%% ============================================================
\subsection{Systematic effects}
\label{sec:systematics}
%% ============================================================

We identify three dominant sources of systematic uncertainty in the
analysis.

\paragraph{Galaxy catalog incompleteness.}
As discussed in Sec.~\ref{sec:pdet_completeness}, the primary source of
systematic bias is the incompleteness of the GLADE+ catalog at
$z \gtrsim 0.03$. Without the completeness correction, the posterior is
systematically pulled toward lower $\HH$ because the catalog is
preferentially populated at low redshifts, where a smaller $h$ can
reproduce the observed $\dL$. The magnitude of this effect is
\pending{Rank-1 budget}. The completeness
correction of Eq.~\eqref{eq:completeness_correction} mitigates this bias,
though its effectiveness depends on the accuracy of the completeness
function $f(z, h)$, which is currently based on digitized data from
Ref.~\cite{Dalya:2021ewn}.

\paragraph{Number of candidate host galaxies.}
Without the MBH mass channel, each GW detection admits $\sim 10$--$10^4$
candidate host galaxies within the $3\sigma$ sky localization and distance
error volume. The true host contribution to the single-event likelihood
can be as small as $\sim 0.6\%$ for high-redshift detections
($z \gtrsim 0.12$), where the galaxy density is high and the distance
measurement is less constraining.
\pending{Rank-1 census of candidates per event with and without the mass
window, GLADE+ today}.

\paragraph{Detection probability normalization.}
The detection probability $\pdet$ enters both the numerator and
denominator of the single-event likelihood. An $h$-dependent bias in the
$\pdet$ estimate can propagate into a systematic shift of the posterior.
In the $\pdet = 1$ baseline, this effect is absent by construction, but
with a realistic $\pdet$ estimated from an injection campaign of
finite size, the normalization carries a statistical uncertainty of
$\sim 1/\sqrt{N_\mathrm{inj}}$ per grid bin.
\pending{Quantify residual $\pdet$ normalization effect from production
run.}

Table~\ref{tab:results_summary} summarizes the key results for both
inference channels.

\begin{table}[t]
  \caption{Summary of $\HH$ constraints from EMRI dark siren inference.
  The injected value is $h_\mathrm{true} = 0.73$.
  \pending{Update with production run values.}}
  \label{tab:results_summary}
  \centering
  \begin{tabular}{l c c}
    \hline\hline
    & Without $\Mz$ & With $\Mz$ \\
    \hline
    MAP estimate &
      \pending{X.XXX} &
      \pending{X.XXX} \\
    $1\sigma$ width &
      \pending{0.0XX} &
      \pending{0.0XX} \\
    Precision (\%) &
      \pending{X.X} &
      \pending{X.X} \\
    Bias $\Delta h / h_\mathrm{true}$ (\%) &
      \pending{X.X} &
      \pending{X.X} \\
    $N_\mathrm{det}$ used &
      \multicolumn{2}{c}{\pending{NNN}} \\
    \hline\hline
  \end{tabular}
\end{table}

%% CITATIONS NEEDED (for gpd-bibliographer):
%% None — all citations used in this section are already in references.bib.
